\documentclass[portuguese]{sbc2025}%
\usepackage{booktabs}
%\usepackage{graphicx}% já está incluido na classe
%\usepackage[utf8]{inputenc} % obsoleto, desnecessário
\usepackage[misc,geometry]{ifsym}
%%%%  %\usepackage{fontspec}
%% Os problemas com a classe eram originados por carregarem ambos  e fontspec simultaneamente. Agora a classe detecta
%% qual engine está em uso e se for o pdflates então carrega o
%% fontenc. Se forem o xelatex ou lualatex então carrega o
%% fontspec. Eu particularemnte prefiro o lualatex ou o xetex
% \usepackage{fontawesome} %%% incluido na classe. Desnecessário
% chamá-lo aqui
%%%% \usepackage{academicons}
%%%% outro encrenqueiro que se tornou desnecessario. Deste typeface só
%%%% se usava o glifo do orcid e o codigo interno bombava.
%%%% Substitui pelo pacote orcidlink (carregado internamete na classe)
%%%% e consertei o codigo na classe.
%\usepackage{color} % carregada dentro da classe
%\usepackage{hyperref} % carregada dentro da classe
\usepackage{aas_macros}
\usepackage[bottom]{footmisc}
%\usepackage{supertabular}
%\usepackage{multicol}
%\usepackage{multirow}
%% O pacote tabularray é muito melhor para
% a criacao de tabela complexas e/ou loooongas; tudo de modo simples.
%% Torna o uso de multicol e multirow desnecessarios
\usepackage{tabularray}
\usepackage{afterpage}
\usepackage{url}
\usepackage{pifont}
\setcitestyle{square}
\definecolor{engtitle}{rgb}{0.5,0.5,0.5}
\definecolor{orcidlogo}{rgb}{0.37,0.48,0.13}
\definecolor{unilogo}{rgb}{0.16, 0.26, 0.58}
\definecolor{maillogo}{rgb}{0.58, 0.16, 0.26}
\definecolor{darkblue}{rgb}{0.0,0.0,0.0}
\hypersetup{colorlinks,breaklinks,
            linkcolor=darkblue,urlcolor=darkblue,
            anchorcolor=darkblue,citecolor=darkblue}
%\hypersetup{colorlinks,citecolor=blue,linkcolor=blue,urlcolor=blue}
%%%%%%% IMPORTANT: We disable hyperlinks by default with this line, to avoid the error "\pdfendlink ended up in different nesting level" while writing.
%\hypersetup{draft}
\category{Artigo de Pesquisa/Research Paper}
\title[Análise Experimental Comparativa de Algoritmos de Criptografia Pós-Quântica do NIST: Desempenho Temporal, Otimização de Compilação e Escalabilidade]{Análise Experimental Comparativa de Algoritmos de Criptografia Pós-Quântica do NIST: Desempenho Temporal, Otimização de Compilação e Escalabilidade}
\engtitle{\textcolor{engtitle}{Comparative Experimental Analysis of NIST Post-Quantum Cryptography Algorithms: Execution Time, Compilation Optimization, and Scalability}}
%THE ORCID IS MANDATORY FOR EACH AUTHOR IN JBCS
\author[Zandona et al. 2025]{
\affil{\textbf{Saulo de Moura Zandona Freitas}~\textcolor{blue}{\faEnvelopeO}~~[{Pontifícia Universidade Católica de Minas Gerais}~|\href{mailto:saulomzf@gmail.com}{~{\textit{saulomzf@gmail.com}}}~]}
\affil{\textbf{Matheus Alcântara Souza~\textcolor{blue}{\faEnvelopeO}}~~[{Pontifícia Universidade Católica de Minas Gerais}~|\href{mailto:matheusalcantara@pucminas.br}{{\textit{matheusalcantara@pucminas.br}}}~]}
}
\begin{document}
\begin{frontmatter}
\maketitle
\begin{mail}
Depto. de Ciência da Computação -- Pontifícia Universidade Católica de Minas Gerais, R. Dom José Gaspar, 500 - Coração Eucarístico, Belo Horizonte - MG, 30535-901
\end{mail}
\begin{abstract-pt}
O avanço dos computadores quânticos representa uma ameaça
significativa aos algoritmos criptográficos clássicos amplamente
utilizados atualmente. Em resposta a esse cenário, a criptografia
pós-quântica tem sido objeto de intensas pesquisas, visando o
desenvolvimento de algoritmos resistentes a ataques quânticos.
Este trabalho apresenta uma análise experimental do desempenho
de algoritmos de criptografia pós-quântica selecionados, com
foco nos finalistas e candidatos do processo de padronização
do NIST.


Os experimentos foram conduzidos utilizando a biblioteca
\textit{liboqs}, considerando mecanismos de encapsulamento de
chaves e esquemas de assinatura digital: ML-KEM (Kyber),
HQC, ML-DSA (Dilithium), FN-DSA (Falcon) e SLH-DSA
(SPHINCS+). Foram avaliadas métricas como tempo de execução
das principais operações criptográficas, tamanho de chaves e
assinaturas, impacto de configurações de compilação,
escalabilidade com o tamanho da mensagem e consumo de
memória (\textit{stack} e \textit{heap}). Os resultados obtidos
permitem uma comparação objetiva entre os algoritmos analisados,
evidenciando diferenças relevantes de desempenho e custo
computacional.


A partir da análise realizada, o trabalho contribui para uma
melhor compreensão dos impactos práticos da adoção da
criptografia pós-quântica, fornecendo subsídios para a escolha
de algoritmos em diferentes cenários de uso e estabelecendo
uma base experimental para estudos futuros.
\end{abstract-pt}
\begin{abstract-en}
The advancement of quantum computing poses a significant threat
to widely used classical cryptographic algorithms. In response
to this scenario, post-quantum cryptography has become an
active research area, aiming to develop algorithms resistant
to quantum attacks. This work presents an experimental
performance analysis of selected post-quantum cryptographic
algorithms, focusing on finalists and candidates from the
NIST standardization process.


The experiments were conducted using the \textit{liboqs}
library, considering key encapsulation mechanisms and digital
signature algorithms: ML-KEM (Kyber), HQC, ML-DSA
(Dilithium), FN-DSA (Falcon), and SLH-DSA (SPHINCS+).
Metrics evaluated include execution time of the main
cryptographic operations, key and signature sizes, the impact
of compilation configurations, scalability with message
size, and memory consumption (stack and heap). The obtained
results enable an objective comparison among the analyzed
algorithms, highlighting relevant differences in performance
and computational cost.


Based on the experimental analysis, this work contributes to
a better understanding of the practical implications of
adopting post-quantum cryptography, providing support for
algorithm selection in different usage scenarios and
establishing an experimental baseline for future studies.
\end{abstract-en}
\begin{pchaves}
Computação Quântica, Criptografia pós-quântica, avaliação de desempenho
\end{pchaves}
\begin{keywords}
Quantum Computing, Post-Quantum Cryptography, performance evaluation
\end{keywords}
%\begin{license}
%Published under the Creative Commons Attribution 4.0 International Public License (CC BY 4.0)
%\end{license}
\end{frontmatter}
\section{Introdução}
\label{sec:intro}
A computação quântica vem se consolidando como um novo paradigma tecnológico, deixando o campo teórico e tornando-se uma área de intenso desenvolvimento. Baseia-se em fenômenos da mecânica quântica, como superposição e emaranhamento, para realizar operações impossíveis à computação clássica. Enquanto os computadores tradicionais utilizam bits, que assumem valores 0 ou 1, os computadores quânticos operam com qubits, capazes de representar ambos os estados simultaneamente até o momento da medição \cite{Livro:nielsen:2010,feynman:82}. Essa propriedade da superposição, aliada à outros fenômenos da quântica, como o emaranhamento, permite correlações entre qubits e processamento paralelo em larga escala, conferindo à computação quântica um poder computacional sem precedentes.
\subsection{Problema}
Com o avanço da computação quântica, essa nova tecnologia representa uma ameaça direta aos sistemas criptográficos clássicos, como RSA e ECC, vulneráveis a ataques quânticos \cite{shor:94}. Mesmo os algoritmos simétricos, como AES e SHA, exigem chaves maiores para manter a segurança \cite{grover:96}. Em resposta, o NIST iniciou o processo de padronização da Criptografia Pós-Quântica (PQC) \cite{NIST:2024}, com o objetivo de estabelecer algoritmos resistentes a ataques quânticos. Contudo, ainda há incertezas sobre o desempenho prático dos finalistas quanto ao tempo de execução
, tamanho das chaves e viabilidade de implementação, o que dificulta sua adoção em larga escala. Surge, assim, a questão central: quais são as diferenças de desempenho entre os algoritmos finalistas do NIST PQC e suas implicações para aplicações reais de segurança da informação?
\subsection{Objetivos}
Com base no processo da padronização do NIST de Criptografia Pós-Quântica, este trabalho busca compreender melhor o desempenho dos algoritmos de criptografia pós-quântica finalistas do processo de padronização do NIST.
\subsubsection{Objetivo principal}
Portanto, o objetivo principal desta monografia é realizar uma análise comparativa de desempenho dos algoritmos finalistas do NIST PQC, avaliando métricas como tempo de execução, consumo de memória, tamanho de chave e eficiência em diferentes cenários de aplicação, a fim de identificar suas potencialidades e limitações para adoção prática em sistemas de segurança da informação.
\subsubsection{Objetivos específicos}
Para alcançar o objetivo principal, estabelecem-se os seguintes objetivos específicos:
\begin{itemize}
    \item Explorar os algoritmos Finalistas do NIST PQC para um melhor entendimento dos mesmos;
    \item buscar um ambiente estável para a implementação os algoritmos;
    \item Definir e aplicar métricas de avaliação de desempenho, a serem analisadas entre os algoritmos;
    \item Comparar e interpretar os resultados obtidos nos experimentos, destacando pontos fortes e limitações de cada algoritmo;
    \item Apontar recomendações sobre a viabilidade prática dos algoritmos em diferentes contextos de uso na segurança da informação.
\end{itemize}
\subsection{Contribuições}
Este trabalho pretende contribuir de forma acadêmica e técnica para a área de segurança da informação ao oferecer uma análise comparativa de desempenho dos algoritmos finalistas do processo de padronização do NIST em criptografia pós-quântica. A partir da realização de experimentos práticos e da avaliação de métricas como tempo de execução, consumo de memória e tamanho de chave, busca-se identificar as principais vantagens e limitações de cada proposta. Além de fornecer subsídios para a escolha de algoritmos mais adequados a diferentes contextos de aplicação, a pesquisa também contribui para ampliar o entendimento da comunidade acadêmica e profissional acerca da viabilidade prática da criptografia pós-quântica.
\subsection{organização da monografia}
Este artigo está organizado da seguinte forma. A Seção~2
descreve os Trabalhos Correlatos, situando esta pesquisa
em relação aos estudos existentes sobre desempenho de
algoritmos PQC. A Seção~3 detalha a Metodologia, definindo
o ambiente de testes, os procedimentos de avaliação e as
métricas adotadas. A Seção~4 descreve a implementação dos
experimentos, incluindo a integração com a biblioteca
\textit{liboqs}, o mecanismo de medição de tempo e a
infraestrutura de perfilamento de memória. A
Seção~5 apresenta os resultados e as análises nas cinco
dimensões avaliadas: tamanho dos artefatos, desempenho
temporal, impacto das configurações de compilação,
escalabilidade com o tamanho da mensagem e consumo de
memória. Por fim, a
Seção~6 apresenta as conclusões do trabalho e aponta
direções para pesquisas futuras.
\section{Trabalhos Correlatos}
\label{sec:Trabalhos}
A transição para a \textit{criptografia pós-quântica} (PQC) representa um dos maiores desafios atuais da segurança digital. Com o avanço dos computadores quânticos, esquemas clássicos como RSA e ECC tornam-se vulneráveis, motivando o processo de padronização conduzido pelo NIST desde 2016. Após cinco rodadas, foram definidos quatro algoritmos principais: \textit{CRYSTALS-Kyber} (KEM), \textit{CRYSTALS-Dilithium} e \textit{Falcon} (assinaturas baseadas em reticulados), e \textit{SPHINCS+} (assinatura baseada em funções hash). Apesar da consolidação teórica, ainda faltam estudos experimentais que avaliem seu desempenho sob diferentes ambientes e configurações de compilação.
\cite{Sosnowski:2023} analisaram a integração de algoritmos PQC ao protocolo TLS~1.3, avaliando KEMs e assinaturas  incluindo \textit{Kyber}, \textit{HQC}, \textit{Dilithium} e \textit{Falcon} quanto à latência e ao tempo de \textit{handshake}. Os resultados mostraram que as soluções pós-quânticas podem atingir desempenho comparável ou superior ao TLS clássico, especialmente em modos híbridos. No entanto, o estudo não aborda \textit{SPHINCS+} nem variações de compilador ou otimização.
Em foco similar, \cite{Opilka:2024} compararam \textit{Dilithium}, \textit{Falcon} e \textit{SPHINCS+} em assinaturas digitais utilizando a biblioteca \textit{liboqs}. As métricas de geração de chaves, assinatura e verificação indicaram que \textit{Dilithium} possui melhor equilíbrio entre velocidade e segurança, enquanto \textit{SPHINCS+} se destaca em robustez, porém com custo computacional elevado. O estudo, contudo, não inclui KEMs nem analisa consumo de memória.
O artigo \cite{Raavi:2025}   reforça esses resultados ao testar as mesmas assinaturas em cenários TLS~1.3, observando reduções de até 68\,\% na latência frente ao RSA-4096, mas sem incorporar esquemas de encapsulamento ou variações de compilação.
\cite{Demir:2025} realizaram uma comparação direta entre o KEM \textit{Kyber} e a assinatura \textit{Dilithium}, medindo tempos de geração de chaves, encapsulamento e verificação em plataformas embarcadas e servidores x86-64. Concluíram que \textit{Kyber} é mais eficiente em encapsulamento e \textit{Dilithium} mantém verificação estável, mas o estudo não inclui \textit{Falcon} e \textit{SPHINCS+}, nem controla \textit{flags} de compilação.
Por fim, \cite{abbasi:2025} apresentaram um \textit{benchmark} multiplataforma abrangendo \textit{Kyber}, \textit{NTRU}, \textit{Dilithium} e \textit{Falcon}, avaliando tempos de operação e níveis de segurança NIST~1,~3 e~5. Apesar do amplo escopo, o trabalho não inclui \textit{SPHINCS+} e carece de análise detalhada sobre otimizações e uso de recursos.
Em síntese, como explicitado na tabela \ref{tab:correlatos}, os estudos revisados avançam na avaliação prática dos algoritmos PQC, mas nenhum realiza uma análise conjunta dos quatro finalistas do NIST sob condições experimentais homogêneas. Assim, este TCC propõe uma avaliação sistemática e reprodutível de \textit{Kyber}, \textit{Dilithium}, \textit{Falcon} e \textit{SPHINCS+}, explorando o impacto de diferentes \textit{flags} de compilação (\texttt{-O0} a \texttt{-O3}, \texttt{-march}, LTO) e níveis de segurança no desempenho, uso de memória e eficiência, contribuindo com uma base comparativa sólida para futuras implementações pós-quânticas.
\begin{table*}[ht]
\caption{Comparação entre trabalhos correlatos e a proposta deste TCC}
\label{tab:correlatos}
\resizebox{\textwidth}{!}{%
\begin{tabular}{p{3.5cm} p{2.1cm} p{3cm} p{2.3cm} p{1.7cm} p{1.8cm} p{2.7cm}}
\toprule
\textbf{Autor / Ano} & \textbf{Ambiente de Teste} & \textbf{Algoritmos Avaliados} & \textbf{Métricas Principais} & \textbf{Variação de Compilação} & \textbf{implementa SPHINCS+?} & \textbf{Foco do Estudo} \\
\midrule
\cite{Sosnowski:2023} & TLS~1.3 (rede) & Kyber, HQC, Dilithium, Falcon & Latência, tempo de handshake & Não & Não & Integração PQC em protocolos TLS \\
\midrule
\cite{Opilka:2024} & Software (liboqs) & Dilithium, Falcon, SPHINCS+ & Geração de chaves, assinatura e verificação & Não & Sim & Desempenho de assinaturas digitais \\
\midrule
\cite{Raavi:2025} & TLS~1.3 (simulado) & Dilithium, Falcon, SPHINCS+ & Tempo e tamanho do handshake & Não & Sim & Assinaturas PQC em TLS \\
\midrule
\cite{Demir:2025} & Embarcado / x86-64 & Kyber, Dilithium & Keygen, encaps, sign, verify & Não & Não & Comparação entre KEM e assinatura \\
\midrule
\cite{abbasi:2025} & Multiplataforma (IoT / Cloud) & Kyber, NTRU, Dilithium, Falcon & Keygen, sign/verify, tempo total & Não & Não & Benchmark de múltiplos ambientes \\
\midrule
\textbf{Este trabalho} & Desktop / Servidor (controlado) &
ML-KEM, HQC, ML-DSA, FN-DSA, SLH-DSA &
Tempo de exec., tamanho de artefatos, variação de flags, níveis NIST e consumo de memória (\textit{stack}/\textit{heap}) &
\textbf{Sim} (\texttt{-O0--O3}, \texttt{-march}, LTO) &
\textbf{Sim} &
Avaliação experimental completa e reprodutível em diferentes níveis de execução \\
\bottomrule
\end{tabular}
} % fecha resizebox
\end{table*}
% ============================================================
% SECAO 3 — Metodologia (foco: O QUE e POR QUE)
% ============================================================
\section{Metodologia}
\label{sec:Metodologia}
A avaliação experimental dos algoritmos pós-quânticos exige um ambiente
de testes rigorosamente definido, de modo a garantir a
reprodutibilidade, precisão e comparabilidade entre os resultados
obtidos. Esta seção descreve os materiais utilizados, os critérios de
seleção dos algoritmos e variantes avaliados, as métricas de desempenho
adotadas e os procedimentos experimentais seguidos. Os detalhes técnicos
da implementação, estrutura do código, integração com a biblioteca e
mecanismo de medição são apresentados na Seção~\ref{sec:Implementacao}.
\subsection{Materiais e Ambiente de Testes}
A caracterização detalhada do ambiente experimental é fundamental para
garantir a reprodutibilidade e a confiabilidade dos resultados obtidos.
\subsubsection{Hardware e Sistema Operacional}
Os experimentos foram realizados em um notebook ASUS TUF Gaming A15
(modelo FA506NCR), equipado com o processador AMD Ryzen 7 7435HS,
pertencente à microarquitetura Rembrandt (6~nm). O processador possui
8 núcleos e 16 threads, operando a até 4,55~GHz em modo turbo, e
oferece suporte a instruções modernas relevantes para algoritmos
pós-quânticos, incluindo x86-64, SSE (1–4.2), SSSE3, AVX, AVX2, FMA3
e SHA. O sistema dispõe de 16~GB de memória DDR5 em modo dual-channel
a 4800~MHz efetivos e armazenamento SSD NVMe.
A placa-mãe utilizada é uma ASUSTeK FA506NCR, versão 1.0, com chipset
AMD Ryzen SOC / AMD FCH, e BIOS American Megatrends International LLC
FA506NCR.304, atualizada em 08/01/2025, executando em modo UEFI com
Secure Boot habilitado. A GPU NVIDIA GeForce RTX 3050 presente no
equipamento não participou da execução dos algoritmos avaliados.
O ambiente de software consiste no Microsoft Windows 11 Home,
versão 10.0.26200 (Build 26200), com configurações de energia ajustadas
para desempenho máximo e processos de segundo plano minimizados. Os
experimentos foram executados sob o Windows Subsystem for Linux (WSL),
com a distribuição Ubuntu 24.04.1 LTS, permitindo o uso nativo de
ferramentas Linux sem virtualização completa.
\subsubsection{Compilador e Dependências}
A compilação dos algoritmos foi realizada com o compilador GCC versão
13.3.0, disponibilizado pela distribuição Ubuntu 24.04.1 LTS sob WSL.
As ferramentas de construção CMake versão 3.28.3 e Ninja-build versão
1.11.1 foram utilizadas nos scripts de compilação da biblioteca e do
projeto experimental.
Para avaliar o impacto das otimizações de compilação no desempenho,
foram empregadas as flags descritas na Tabela~\ref{tab:flags}. Essa
análise permite quantificar o ganho obtido com vetorização automática,
uso de instruções específicas da arquitetura e otimização em tempo de
linkedição, aspectos relevantes para algoritmos PQC que fazem uso
intensivo de operações aritméticas como NTT e FFT.
\begin{table*}
\centering
\caption{Flags de otimização avaliadas e seus efeitos no desempenho.}
\label{tab:flags}
\begin{tabular}{|c|p{11cm}|}
\hline
\textbf{Flag} & \textbf{Descrição Técnica e Efeitos Esperados} \\
\hline
\texttt{-O0} & Desativa otimizações. Mantém o código literalmente como
escrito, preservando estrutura e instruções redundantes. Serve como
\textit{baseline} para avaliar o custo algorítmico puro. \\
\hline
\texttt{-O1} & Realiza otimizações leves, removendo código morto e
redundâncias simples. Pouco impacto computacional, mas melhora tempos
em rotinas pequenas. \\
\hline
\texttt{-O2} & Ativa otimizações amplas: desenrolamento moderado de
loops, inlining seletivo, propagação de constantes, eliminação de
subexpressões, melhor alocação de registradores. Fornece bom equilíbrio
entre desempenho e tamanho do binário. \\
\hline
\texttt{-O3} & Nível agressivo de otimização. Inclui vetorização
automática (SSE/AVX/AVX2), desenrolamento profundo de loops,
\textit{pipelining} e heurísticas avançadas de reorganização de código.
Produz ganhos substanciais em algoritmos de PQC baseados em NTT, FFT,
hashing intensivo e amostragem gaussiana. \\
\hline
\texttt{-march=native} & Permite ao compilador utilizar instruções
específicas do processador físico (AMD Ryzen 7 7435HS), incluindo AVX2,
BMI2, SHA e FMA3. Essencial para explorar vetorização e aceleração por
hardware. \\
\hline
\texttt{-flto} & Ativa otimização em tempo de linkedição (LTO),
permitindo ao compilador otimizar o programa como um todo, eliminando
fronteiras artificiais entre módulos. Reduz overhead em rotinas
fortemente modularizadas, como as da \texttt{liboqs}. \\
\hline
\end{tabular}
\end{table*}
\subsubsection{Biblioteca Criptográfica}
Para a realização dos experimentos empregou-se a \textit{Open Quantum
Safe Library} (\textbf{liboqs}), mantida pelo projeto Open Quantum Safe
(OQS). A liboqs reúne implementações atualizadas dos esquemas
selecionados pelo NIST, incluindo os algoritmos avaliados neste estudo:
os KEMs ML-KEM (Kyber) e HQC, e os esquemas de assinatura ML-DSA
(Dilithium), FN-DSA (Falcon) e SLH-DSA (SPHINCS+).
A principal vantagem da liboqs é sua API unificada, que permite executar
geração de chaves, encapsulamento, decapsulamento, assinatura e
verificação por meio de uma interface uniforme para todos os algoritmos.
Essa padronização elimina discrepâncias entre os códigos de referência
dos autores, garantindo consistência metodológica na comparação
experimental. A biblioteca disponibiliza múltiplas variantes de cada
algoritmo, cobrindo os níveis de segurança NIST~1, NIST~3 e NIST~5.
\subsection{Procedimentos de Avaliação}
Para cada algoritmo foram avaliadas as operações fundamentais:
geração de chaves, encapsulamento e decapsulamento (KEMs), e
geração de chaves, assinatura e verificação (esquemas de
assinatura). Todas as variantes parametrizadas foram incluídas:
ML-KEM-512, ML-KEM-768 e ML-KEM-1024; ML-DSA-44, ML-DSA-65
e ML-DSA-87; FN-DSA-512 e FN-DSA-1024;
SLH-DSA-SHAKE-128s/f, SLH-DSA-SHAKE-192s/f e
SLH-DSA-SHAKE-256s/f; e HQC-128, HQC-192 e HQC-256,
permitindo comparação entre níveis de segurança dentro de
cada família criptográfica.


Cada operação foi executada 200 vezes, precedidas de 10
iterações de aquecimento para estabilização dos caches do
processador. Esse número de repetições foi definido para
garantir estabilidade estatística das medições, conforme
prática adotada em estudos experimentais de criptografia
pós-quântica~\cite{Opilka:2024, abbasi:2025}.


Além dos testes individuais, foram conduzidos experimentos de
escalabilidade para os algoritmos de assinatura digital,
baseados no aumento progressivo do tamanho da mensagem de
entrada. Os tamanhos avaliados variam de 1~B a 1~MB em
potências de dois (21 pontos), permitindo analisar o impacto
do hashing interno sobre o desempenho de cada família
criptográfica. Essa análise não se aplica aos KEMs, pois suas
operações não recebem mensagem como entrada: o encapsulamento
opera exclusivamente sobre a chave pública, com saídas de
tamanho fixo.


\subsection{Métricas e Critérios de Comparação}


A avaliação foi estruturada em cinco dimensões principais:


\begin{enumerate}
    \item \textbf{Tempo de execução}: métrica primária,
    calculada como média, desvio-padrão, mediana e coeficiente
    de variação (CV\%) das 200 repetições. A mediana é
    preferível à média para comparações, por ser robusta a
    outliers causados por interferências do sistema
    operacional. O CV\% indica a estabilidade do algoritmo:
    valores elevados, como os observados no Sign do ML-DSA,
    refletem o comportamento probabilístico inerente ao
    mecanismo de amostragem por rejeição.


    \item \textbf{Tamanho dos artefatos}: chave pública,
    chave secreta, ciphertext (KEMs) e assinatura (esquemas
    de assinatura), extraídos das estruturas internas da
    liboqs conforme as especificações oficiais de cada esquema.


    \item \textbf{Impacto das flags de compilação}: variação
    no tempo de execução obtida com cada configuração em
    relação ao baseline \texttt{-O0}, avaliando a
    sensibilidade de cada algoritmo à infraestrutura de
    compilação moderna.


    \item \textbf{Escalabilidade}: variação do tempo de
    assinatura e verificação em função do tamanho da mensagem,
    revelando o custo do hashing interno de cada família
    criptográfica.


    \item \textbf{Consumo de memória}: uso de pilha
    (\textit{stack}) e pico de \textit{heap} de cada operação,
    medido por técnicas complementares. O consumo de
    \textit{stack} é obtido via \textit{stack painting} com
    \textit{buffer} sentinela de 4\,MB em \textit{thread}
    POSIX dedicada; o pico de \textit{heap} é obtido com
    \texttt{valgrind --tool=massif}, capturando o valor máximo
    de \texttt{mem\_heap\_B} durante a execução isolada de
    cada operação. Ambas as medições são realizadas com
    compilação \texttt{-O0} para preservar a estrutura real
    de alocações.
\end{enumerate}


As comparações são realizadas entre algoritmos do mesmo nível
de segurança NIST, garantindo equidade na análise.
% ============================================================
% SECAO 4 — Descricao das Implementacoes (foco: COMO)
% ============================================================
\section{Descrição das Implementações}
\label{sec:Implementacao}
Esta seção descreve os aspectos técnicos da implementação experimental:
a organização do código, a integração com a liboqs, o mecanismo de
medição de tempo e o registro dos resultados. Enquanto a
Seção~\ref{sec:Metodologia} define \textit{o que} foi medido e
\textit{por que}, esta seção detalha \textit{como} os experimentos
foram executados, possibilitando a reprodução independente dos
resultados.
\subsection{Estrutura do Projeto}
A implementação foi desenvolvida em linguagem C e organizada de forma
modular, separando claramente os componentes de execução criptográfica,
medição de desempenho e coleta de resultados. A estrutura de diretórios
divide o projeto em módulos independentes para cada família de
algoritmos (KEM e assinatura) e para os utilitários de temporização e
estatística, facilitando a extensão futura com novos algoritmos ou
métricas.
\subsection{Integração com a biblioteca liboqs}
Todas as operações criptográficas foram executadas via liboqs. Para os
KEMs (ML-KEM e HQC), utilizou-se a interface \texttt{OQS\_KEM}, cujo
fluxo de execução segue três etapas:
\begin{enumerate}
    \item Inicialização: \texttt{OQS\_KEM\_new()};
    \item Geração de chaves: \texttt{OQS\_KEM\_keypair()};
    \item Encapsulamento e decapsulamento: \texttt{OQS\_KEM\_encaps()}
          e \texttt{OQS\_KEM\_decaps()}.
\end{enumerate}
Para os algoritmos de assinatura (ML-DSA, FN-DSA e SLH-DSA),
utilizou-se a interface \texttt{OQS\_SIG}:
\begin{enumerate}
    \item Inicialização: \texttt{OQS\_SIG\_new()};
    \item Geração de chaves: \texttt{OQS\_SIG\_keypair()};
    \item Assinatura: \texttt{OQS\_SIG\_sign()};
    \item Verificação: \texttt{OQS\_SIG\_verify()}.
\end{enumerate}
A interface uniforme entre todos os algoritmos reduz a complexidade da
implementação e garante que diferenças nos resultados reflitam
características dos algoritmos, não da infraestrutura de execução.
Ao encerrar cada benchmark, a chave secreta é zerada via
\texttt{OQS\_MEM\_secure\_free()}, seguindo boas práticas de segurança
criptográfica.
\subsection{Implementação dos Benchmarks}
Cada rotina de benchmark executa o seguinte fluxo:
\begin{enumerate}
    \item Alocação dos buffers de chaves, ciphertext e assinatura;
    \item Execução das iterações de aquecimento (sem medição);
    \item Laço principal de 200 iterações, com medição individual
          de cada operação;
    \item Verificação de corretude a cada iteração para KEMs,
          confirmando que os \textit{shared secrets} gerados por
          encapsulamento e decapsulamento são idênticos; para
          assinaturas, confirmando que a verificação retorna sucesso;
    \item Cálculo das estatísticas e escrita no arquivo CSV.
\end{enumerate}
A verificação de corretude a cada iteração é um diferencial
importante: garante que os tempos medidos correspondem a execuções
válidas dos algoritmos, não a falhas silenciosas.
\subsection{Medição do Tempo de Execução}
A medição de tempo utiliza a função \texttt{clock\_gettime} com o
relógio \texttt{CLOCK\_MONOTONIC}, pertencente à API POSIX do Linux.
O procedimento em cada iteração é:
\begin{enumerate}
    \item Leitura do contador imediatamente antes da operação;
    \item Execução da operação criptográfica;
    \item Leitura do contador imediatamente após;
    \item Cálculo da diferença, convertida para microssegundos.
\end{enumerate}
As funções de temporização são implementadas como \texttt{static inline}
no código, eliminando o overhead de chamada de função nas medições
críticas. A escolha do \texttt{CLOCK\_MONOTONIC} garante leituras
estritamente crescentes, imunes a ajustes de NTP ou alterações manuais
do relógio do sistema.
\subsection{Medição do Consumo de Memória}
\label{sec:impl-memoria}


A medição de memória é realizada por dois programas distintos,
compilados com \texttt{-O0 -g} para preservar o comportamento
de alocação sem otimizações do compilador.


\textbf{Stack painting (\texttt{memory\_benchmark.c}).}
Para cada operação (\textit{keygen}, \textit{encap}/\textit{sign},
\textit{decap}/\textit{verify}), uma \textit{thread} POSIX é criada
com \texttt{pthread\_attr\_setstack()} apontando para um
\textit{buffer} de 4\,MB previamente pintado com o byte
sentinela~\texttt{0xAB}.
Após o retorno da operação, o \textit{buffer} é varrido do endereço
mais baixo em direção ao mais alto; o primeiro byte diferente
de~\texttt{0xAB} delimita a região de \textit{stack} não utilizada,
e o uso real é calculado como o complemento.
O resultado reportado é a mediana de cinco execuções após dois
aquecimentos.


\textbf{Valgrind massif (\texttt{massif\_single.c}).}
Um programa auxiliar executa exatamente \emph{uma} chamada da
operação desejada por invocação, permitindo que o
\texttt{valgrind --tool=massif} perfile o \textit{heap} daquela
operação isolada.
O \textit{script} \texttt{run\_massif\_all.sh} automatiza a execução
para todos os algoritmos e operações; o \textit{script}
\texttt{parse\_massif.py} extrai o valor de pico de
\texttt{mem\_heap\_B} de cada arquivo \texttt{.massif} gerado e
consolida os resultados em CSV\@.


Como a versão do \texttt{liboqs} utilizada ainda adota os nomes
internos originais (p.\,ex., \texttt{Falcon-512} em vez de
\texttt{FN-DSA-512}), o programa \texttt{massif\_single.c} inclui
uma tabela de mapeamento que traduz os nomes FIPS/NIST para os
identificadores aceitos pela API\@.


\subsection{Registro e Organização dos Resultados}
Os resultados são exportados em arquivos CSV separados por categoria
(KEMs e assinaturas), com uma linha por algoritmo contendo: nome do
algoritmo, nível NIST, tamanhos dos artefatos, e para cada operação
(KeyGen, Encaps/Sign, Decaps/Verify) os valores de média, desvio-padrão,
mínimo e mediana em microssegundos. Esse formato permite importação
direta em ferramentas de análise estatística e geração de gráficos para
a seção de resultados.
% ============================================================
% SEÇÃO 5 — Implementações, Testes, Resultados e Análises
% ============================================================
\section{Implementações, Testes, Resultados e Análises}
Esta seção apresenta os resultados obtidos nos experimentos
realizados com os algoritmos de criptografia pós-quântica
selecionados, organizados em cinco dimensões de análise.
A primeira examina os tamanhos dos artefatos criptográficos
gerados por cada algoritmo, por ser uma propriedade
estrutural que independe do ambiente de execução.
A segunda avalia o desempenho temporal das operações
fundamentais de cada família, adotando a configuração
\texttt{-O0} como referência de custo algorítmico puro.
A terceira quantifica o impacto das 16 configurações de
compilação avaliadas sobre o tempo de execução, por meio
de mapas de calor que permitem identificar os padrões de
sensibilidade de cada família às otimizações do compilador.
A quarta analisa como os tempos de assinatura e verificação
variam com o tamanho da mensagem, revelando o comportamento
de escalabilidade de cada esquema.
A quinta avalia o consumo de memória das operações
criptográficas por meio de duas técnicas complementares:
\textit{stack painting} para medir o uso de pilha e
\texttt{valgrind --tool=massif} para medir o pico de
\textit{heap}.
As comparações entre algoritmos de famílias distintas são
realizadas dentro do mesmo nível de segurança NIST,
garantindo equidade na análise.
% ------------------------------------------------------------
% 5.1 — Tamanho dos Artefatos Criptográficos
% ------------------------------------------------------------
\subsection{Tamanho dos Artefatos Criptográficos}
O tamanho dos artefatos criptográficos, compreendendo chaves,
ciphertexts e assinaturas, é uma métrica estrutural de cada
algoritmo, determinada por sua especificação e independente
do ambiente de execução ou das configurações de compilação.
A Tabela~\ref{tab:kem-tamanho} reúne os tamanhos dos
artefatos dos algoritmos KEM, e a
Tabela~\ref{tab:sig-tamanho} apresenta os dos esquemas de
assinatura digital.
\begin{table}[!htbp]
\caption{Tamanho dos artefatos dos algoritmos KEM (em bytes).}
\label{tab:kem-tamanho}
\centering\small
\resizebox{\columnwidth}{!}{%
\begin{tabular}{lcrrrr}
\toprule
\textbf{Algoritmo} & \textbf{Nível} &
  \textbf{PK (B)} & \textbf{SK (B)} &
  \textbf{CT (B)} & \textbf{SS (B)} \\
\midrule
ML-KEM-512  & 1 &    800 &  1.632 &    768 & 32 \\
ML-KEM-768  & 3 &  1.184 &  2.400 &  1.088 & 32 \\
ML-KEM-1024 & 5 &  1.568 &  3.168 &  1.568 & 32 \\
\midrule
HQC-128  & 1 &  2.249 &  2.305 &  4.433 & 64 \\
HQC-192  & 3 &  4.522 &  4.586 &  8.978 & 64 \\
HQC-256  & 5 &  7.245 &  7.317 & 14.421 & 64 \\
\bottomrule
\end{tabular}%
}
\end{table}
\begin{table}[!htbp]
\caption{Tamanho dos artefatos dos esquemas de assinatura digital (em bytes).}
\label{tab:sig-tamanho}
\centering\small
\resizebox{\columnwidth}{!}{%
\begin{tabular}{lcrrr}
\toprule
\textbf{Algoritmo} & \textbf{Nível} &
  \textbf{PK (B)} & \textbf{SK (B)} &
  \textbf{Assinatura (B)} \\
\midrule
ML-DSA-44 & 2 &  1.312 &  2.560 &  2.420 \\
ML-DSA-65 & 3 &  1.952 &  4.032 &  3.309 \\
ML-DSA-87 & 5 &  2.592 &  4.896 &  4.627 \\
\midrule
FN-DSA-512  & 1 &   897 &  1.281 &   656\textsuperscript{a} \\
FN-DSA-1024 & 5 & 1.793 &  2.305 & 1.273\textsuperscript{a} \\
\midrule
SLH-DSA-SHAKE-128\textit{s} & 1 &  32 &   64 &  7.856 \\
SLH-DSA-SHAKE-128\textit{f} & 1 &  32 &   64 & 17.088 \\
SLH-DSA-SHAKE-192\textit{s} & 3 &  48 &   96 & 16.224 \\
SLH-DSA-SHAKE-192\textit{f} & 3 &  48 &   96 & 35.664 \\
SLH-DSA-SHAKE-256\textit{s} & 5 &  64 &  128 & 29.792 \\
SLH-DSA-SHAKE-256\textit{f} & 5 &  64 &  128 & 49.856 \\
\bottomrule
\multicolumn{5}{l}{\footnotesize\textsuperscript{a}~Tamanho médio
  observado; a assinatura do FN-DSA é de comprimento variável.}
\end{tabular}%
}
\end{table}
No grupo dos KEMs, o ML-KEM produz artefatos significativamente
mais compactos que o HQC em todos os níveis de segurança.
No nível 5, o ciphertext do HQC-256 alcança 14.421~B, valor
aproximadamente 9,2 vezes superior ao do ML-KEM-1024
(1.568~B). Essa disparidade é consequência das distintas
bases matemáticas adotadas: o ML-KEM opera sobre módulos de
reticulados com estrutura algébrica que viabiliza
representações compactas, ao passo que o HQC fundamenta-se
em códigos corretores de erros, exigindo ciphertexts
proporcionalmente maiores para atingir o mesmo nível de
segurança.
No grupo dos esquemas de assinatura, o FN-DSA produz as
assinaturas mais compactas, com 656~B no nível 1, resultado
da codificação comprimida que explora a estrutura NTRU.
No extremo oposto, as assinaturas do SLH-DSA variam de
7.856~B (variante \textit{s}, nível 1) a 49.856~B (variante
\textit{f}, nível 5), sendo as maiores do conjunto avaliado.
Por outro lado, o SLH-DSA possui as menores chaves públicas
(32 a 64~B), propriedade direta de sua construção baseada
exclusivamente em funções de hash, que dispensa qualquer
estrutura algébrica adicional. O ML-DSA ocupa uma posição
intermediária equilibrada, com chaves e assinaturas moderadas
e crescimento proporcional ao nível de segurança.
Uma análise complementar relevante é o \textit{overhead total}
de armazenamento por algoritmo, definido como a soma de chave
pública, chave secreta e artefato gerado. Esse valor representa
o custo mínimo de memória para operar um par de chaves em
uma única transação. No grupo dos KEMs, o ML-KEM-512 impõe
o menor overhead total (3.200~B), enquanto o HQC-256 exige
28.983~B, diferença de 9,1$\times$ no mesmo nível de
segurança. No grupo das assinaturas, o FN-DSA-512 é o mais
econômico (2.834~B no total), beneficiado por suas assinaturas
comprimidas. No extremo oposto, o SLH-DSA-SHAKE-128f acumula
17.184~B, dos quais 99,4\% se concentram na assinatura.
Uma propriedade estrutural do SLH-DSA merece atenção
particular: a distribuição do overhead é radicalmente
assimétrica. As chaves públicas e privadas somam no máximo
192~B (nível 5), enquanto a assinatura pode superar 49~KiB.
Em contextos onde a chave é armazenada localmente uma única
vez e a assinatura é transmitida repetidamente, como em
infraestruturas de chaves públicas, o custo dominante é a
largura de banda por assinatura, não o armazenamento de
chaves. O FN-DSA inverte essa lógica: assinaturas compactas
(656~B) a um custo elevado de geração de chave, o que o
favorece em cenários de chave de longa duração com muitas
assinaturas por par de chaves.
% ------------------------------------------------------------
% 5.2 — Desempenho Temporal dos Algoritmos
% ------------------------------------------------------------
\subsection{Desempenho Temporal dos Algoritmos}
Os tempos de execução foram medidos com a configuração
\texttt{-O0}, adotada como referência por representar o
custo algorítmico puro, sem interferência de otimizações
do compilador. Cada operação foi executada 200 vezes,
precedidas de 10 iterações de aquecimento. São reportados
a média, o desvio-padrão e a mediana em microssegundos
($\mu$s). A mediana é preferível à média para comparações
por ser robusta a valores atípicos causados por
interferências do sistema operacional.
\subsubsection{Algoritmos KEM}
A Tabela~\ref{tab:kem-tempo} reúne os tempos médios de
execução dos algoritmos de encapsulamento de chaves.
\begin{table}[!htbp]
\caption{Desempenho temporal dos KEMs com \texttt{-O0}
($\mu$s, $n=200$).}
\label{tab:kem-tempo}
\centering\small
\resizebox{\columnwidth}{!}{%
\begin{tabular}{lcrrr}
\toprule
\textbf{Algoritmo} & \textbf{Nível} &
  \textbf{KeyGen} & \textbf{Encaps} & \textbf{Decaps} \\
\midrule
ML-KEM-512  & 1 &   6,77 $\pm$ 1,29 &   7,16 $\pm$ 1,55 &   6,19 $\pm$ 0,16 \\
ML-KEM-768  & 3 &  11,08 $\pm$ 3,45 &  11,10 $\pm$ 3,56 &  10,27 $\pm$ 1,48 \\
ML-KEM-1024 & 5 &  15,01 $\pm$ 1,33 &  15,42 $\pm$ 1,43 &  14,85 $\pm$ 2,65 \\
\midrule
HQC-128  & 1 & 1.398 $\pm$ 573 &  2.712 $\pm$  17 &  4.114 $\pm$  55 \\
HQC-192  & 3 & 4.161 $\pm$  45 &  8.353 $\pm$ 176 & 12.593 $\pm$ 221 \\
HQC-256  & 5 & 7.403 $\pm$ 185 & 14.806 $\pm$ 365 & 22.322 $\pm$ 232 \\
\bottomrule
\end{tabular}%
}
\end{table}
O ML-KEM supera o HQC em todas as operações e em todos os
níveis de segurança avaliados. No nível 1, o ML-KEM-512
executa o KeyGen em 6,77~$\mu$s, enquanto o HQC-128 requer
1.398~$\mu$s para a mesma operação, razão de aproximadamente
206$\times$. A diferença é mais acentuada no Decaps: o
HQC-128 consome 4.114~$\mu$s contra 6,19~$\mu$s do
ML-KEM-512, fator de 665$\times$. No nível 5, o HQC-256
registra Decaps de 22.322~$\mu$s, valor 1.503$\times$ superior
aos 14,85~$\mu$s do ML-KEM-1024.
Essa vantagem resulta de dois fatores combinados: a estrutura
algébrica baseada em módulos de reticulados viabiliza
operações mais eficientes por meio da Transformada
Número-Teórica (NTT), e a liboqs fornece implementação com
instruções AVX2 nativas do processador AMD Ryzen 7 7435HS.
O HQC, fundamentado em teoria de códigos, não dispõe de
aceleração equivalente na versão avaliada.
O desvio-padrão elevado do HQC-128 no KeyGen (573~$\mu$s,
CV~$\approx$~41\%) indica sensibilidade às variações de
carga do sistema nessa operação específica, ao passo que
Encaps e Decaps mantêm comportamento estável ao longo das
200 iterações.
A análise do crescimento entre níveis de segurança revela
comportamentos distintos entre as famílias. O ML-KEM exibe
crescimento aproximadamente linear: o KeyGen do ML-KEM-1024
consome 2,2$\times$ mais tempo que o do ML-KEM-512, razão
próxima da esperada pelo aumento do parâmetro modular. O HQC
apresenta crescimento supralinear: o KeyGen do HQC-256 consome
5,3$\times$ mais que o do HQC-128, reflexo do custo de
codificação e decodificação em esquemas baseados em códigos,
cuja complexidade cresce mais que linearmente com o parâmetro
de segurança.
A assimetria entre operações também difere entre as famílias.
No ML-KEM, as três operações são equivalentes em custo: Decaps
representa apenas 0,86$\times$ a 0,99$\times$ o tempo de
Encaps, e nenhuma operação excede duas vezes o custo das
demais. No HQC, o Decaps é sistematicamente 1,5$\times$ mais
lento que o Encaps e 3,0$\times$ mais lento que o KeyGen em
todos os níveis. Esse padrão é intrínseco ao processo de
decodificação de síndromes em códigos corretores de erros,
que exige mais operações que a codificação correspondente.
\subsubsection{Esquemas de Assinatura Digital: ML-DSA e FN-DSA}
A Tabela~\ref{tab:sig-lattice-tempo} reúne os tempos de
execução dos esquemas de assinatura baseados em reticulados.
O coeficiente de variação (CV\%) da operação de assinatura
é incluído por revelar diferenças estruturais relevantes
entre os dois algoritmos.
\begin{table}[!htbp]
\caption{Desempenho de ML-DSA e FN-DSA com \texttt{-O0}
($\mu$s, $n=200$).}
\label{tab:sig-lattice-tempo}
\centering\small
\resizebox{\columnwidth}{!}{%
\begin{tabular}{lcrrrr}
\toprule
\textbf{Algoritmo} & \textbf{Nível} &
  \textbf{KeyGen} & \textbf{Sign} & \textbf{Verify} &
  \textbf{CV\textsubscript{sign}} \\
\midrule
ML-DSA-44 & 2 &  22,25 &  57,42 & 21,35 & 62,0\% \\
ML-DSA-65 & 3 &  40,56 &  96,71 & 38,14 & 58,4\% \\
ML-DSA-87 & 5 &  65,05 & 113,72 & 61,35 & 51,5\% \\
\midrule
FN-DSA-512  & 1 & 4.913,81 & 175,32 & 34,39 & 3,4\% \\
FN-DSA-1024 & 5 & 14.772,44 & 352,29 & 68,28 & 2,2\% \\
\bottomrule
\end{tabular}%
}
\end{table}
O ML-DSA mantém KeyGen, Sign e Verify na mesma ordem de
grandeza, o que representa o melhor equilíbrio entre
operações dentre os esquemas avaliados. O CV elevado do
Sign (62,0\% no ML-DSA-44) é consequência do mecanismo de
\textit{Fiat-Shamir with Aborts}: o algoritmo rejeita
candidatos de assinatura que não satisfazem os critérios
de segurança e reinicia a amostragem, gerando uma
distribuição assimétrica com cauda longa à direita. Nesses
casos, a mediana (44,98~$\mu$s no ML-DSA-44) é
estatisticamente mais representativa do tempo típico que a
média (57,42~$\mu$s).
A razão entre a média e a mediana do Sign fornece uma medida
direta do viés introduzido pela amostragem por rejeição. No
ML-DSA-44, a média supera a mediana em 27,7\%; no ML-DSA-65
essa diferença cai para 23,4\%; e no ML-DSA-87 para 18,1\%.
Esse decréscimo é interpretável: variantes de nível de
segurança mais alto utilizam parâmetros que reduzem a taxa
de rejeição do protocolo, tornando a distribuição de tempos
progressivamente menos assimétrica. A implicação prática é
que sistemas que dimensionam timeouts com base na média do
ML-DSA-44 podem subestimar o tempo real de execução em até
28\% para um conjunto típico de requisições.
O FN-DSA exibe perfil assimétrico: o KeyGen consome
4.913,81~$\mu$s no nível 1, valor 28$\times$ superior ao
Sign (175,32~$\mu$s) e 143$\times$ superior ao Verify
(34,39~$\mu$s). O CV do Sign, inferior a 4\%, confirma
comportamento determinístico nessa fase. O custo elevado
do KeyGen é intrínseco ao processo de amostragem da base
NTRU sobre a rede \textit{fast Fourier}, não constituindo
anomalia experimental. Em cenários de chave efêmera, como
protocolos que geram um novo par a cada sessão, esse custo
é pago integralmente a cada conexão. Quando o par de chaves
é reutilizado ao longo de múltiplas assinaturas, o FN-DSA
torna-se competitivo: a partir da segunda assinatura, o
custo amortizado por operação aproxima-se do Sign isolado
(175~$\mu$s), desempenho inferior apenas ao ML-DSA-44
(57~$\mu$s) entre os esquemas de nível 1.
\subsubsection{Esquemas de Assinatura Digital: SLH-DSA}
A Tabela~\ref{tab:slh-tempo} apresenta os tempos do SLH-DSA
em separado, pois as operações de KeyGen e Sign operam na
escala de milissegundos, incompatível com a escala de
microssegundos dos demais esquemas.
\begin{table}[!htbp]
\caption{Desempenho do SLH-DSA com \texttt{-O0} (KeyGen
e Sign em ms; Verify em $\mu$s; $n=200$).}
\label{tab:slh-tempo}
\centering\small
\resizebox{\columnwidth}{!}{%
\begin{tabular}{lcrrr}
\toprule
\textbf{Algoritmo} & \textbf{Nível} &
  \textbf{KeyGen (ms)} & \textbf{Sign (ms)} &
  \textbf{Verify ($\mu$s)} \\
\midrule
SLH-DSA-SHAKE-128\textit{s} & 1 &  34,22 & 260,40 &  330,20 \\
SLH-DSA-SHAKE-128\textit{f} & 1 &   0,57 &  13,39 &  896,51 \\
\midrule
SLH-DSA-SHAKE-192\textit{s} & 3 &  50,87 & 460,57 &  477,00 \\
SLH-DSA-SHAKE-192\textit{f} & 3 &   0,83 &  21,40 & 1.345,31 \\
\midrule
SLH-DSA-SHAKE-256\textit{s} & 5 &  34,76 & 413,40 &  691,31 \\
SLH-DSA-SHAKE-256\textit{f} & 5 &   2,18 &  43,41 & 1.309,24 \\
\bottomrule
\end{tabular}%
}
\end{table}
O SLH-DSA registra o maior custo de assinatura dentre todos
os esquemas avaliados. A variante \textit{s}, que prioriza
assinaturas menores, requer até 460,57~ms por operação
(SLH-DSA-SHAKE-192s), tornando-a inadequada para aplicações
com alta frequência de assinaturas. A variante \textit{f}
reduz esse custo entre 19$\times$ e 22$\times$, com a
contrapartida de assinaturas maiores, conforme a
Tabela~\ref{tab:sig-tamanho}.
O comportamento se inverte na verificação: a variante
\textit{f} é até 2$\times$ mais lenta que a \textit{s} no
Verify, pois percorre um número maior de nós nas árvores de
Merkle da estrutura FORS. Esse padrão traduz o compromisso
fundamental do SLH-DSA: assinaturas menores implicam maior
custo de geração, e menor custo de geração implica
assinaturas maiores com verificação mais custosa.
A comparação quantitativa entre as variantes revela que o
trade-off não é uniforme entre os níveis de segurança. Para
os níveis 1 e 3, o Sign da variante \textit{s} é
aproximadamente 20$\times$ mais lento que o da \textit{f}
(19,4$\times$ e 21,5$\times$, respectivamente), e o KeyGen
é cerca de 60$\times$ mais lento (59,8$\times$ e
61,0$\times$). No nível 5, essa disparidade cai para
9,5$\times$ no Sign e 15,9$\times$ no KeyGen, indicando
que os parâmetros da variante \textit{s} de 256 bits foram
ajustados para uma relação mais equilibrada entre tamanho de
assinatura e custo computacional. A verificação segue padrão
análogo: a variante \textit{f} é 2,72$\times$, 2,82$\times$
e 1,89$\times$ mais lenta no Verify que a \textit{s}, para
os níveis 1, 3 e 5.
O CV do Sign do SLH-DSA é inferior a 2\% em todas as
variantes, confirmando comportamento determinístico em
contraste direto com o ML-DSA, cujo CV ultrapassa 50\%.
% ------------------------------------------------------------
% 5.3 — Impacto das Configurações de Compilação
% ------------------------------------------------------------
\subsection{Impacto das Configurações de Compilação}
Esta seção quantifica o efeito das 16 configurações de
compilação sobre o desempenho de cada família criptográfica.
As configurações resultam da combinação dos níveis
\texttt{-O0} a \texttt{-O3} com as flags auxiliares
\texttt{-march=native} e \texttt{-flto}, aplicadas
isoladamente e em conjunto. As
Figuras~\ref{fig:heatmap-kem}, \ref{fig:heatmap-sig} e
\ref{fig:heatmap-slh} apresentam mapas de calor com os
tempos de execução reais anotados em cada célula,
normalizados por linha para destacar a variação relativa
de cada algoritmo independentemente de sua magnitude
absoluta. Células verdes identificam as configurações mais
rápidas; vermelhas, as mais lentas.
\subsubsection{Algoritmos KEM}
\begin{figure}[!htbp]
  \centering
  \includegraphics[width=\linewidth]{graficos/g6a_heatmap_kem.pdf}
  \caption{Mapa de calor do tempo de Decaps dos KEMs por
  configuração de compilação. Verde indica a configuração
  mais rápida; vermelho, a mais lenta. Valores em $\mu$s
  (ML-KEM) e ms (HQC), normalizados por linha.}
  \label{fig:heatmap-kem}
\end{figure}
O mapa de calor da Figura~\ref{fig:heatmap-kem} revela um
padrão inesperado para o ML-KEM: a configuração \texttt{-O0}
figura entre as mais rápidas para todas as três variantes,
enquanto \texttt{-O2} e \texttt{-O2~-flto} concentram as
células mais vermelhas. O ML-KEM-512, por exemplo, registra
Decaps de 6,2~$\mu$s com \texttt{-O0} e 6,9~$\mu$s com
\texttt{-O2~-flto}, regressão de 11\%. Esse padrão decorre
da arquitetura interna da liboqs: as implementações do
ML-KEM utilizam código AVX2 escrito em assembly, compilado
com otimizações próprias da biblioteca, independentemente
das flags aplicadas ao código do \textit{harness} de
benchmark. As variações observadas refletem diferenças no
código de temporização e na pressão sobre o cache de
instruções introduzida pelo binário do \textit{harness},
não ganhos reais no núcleo criptográfico. Para o ML-KEM,
a escolha das flags de compilação do código cliente é
irrelevante para o desempenho final. Além disso, as
flutuações observadas chegam a 0,89$\times$ no ML-KEM-512
com \texttt{-O2\_lto}, magnitude comparável ao desvio-padrão
das medições individuais, o que impede atribuir as variações
a qualquer efeito real das flags.
O HQC exibe padrão análogo: \texttt{-O0} e \texttt{-O3}
alternam entre as células mais verdes, enquanto
\texttt{-O2~-flto} e \texttt{-O2~-march=native} produzem
os piores tempos. O HQC-256 Decaps varia de 22,3~ms com
\texttt{-O0} a 24,7~ms com \texttt{-O2~-flto}, diferença
de 10,7\%. Tal resultado indica que as implementações do
HQC na liboqs já incorporam otimizações internas que não
são amplificadas pelas flags de compilação aplicadas ao
\textit{harness}.
\subsubsection{ML-DSA e FN-DSA}
\begin{figure}[!htbp]
  \centering
  \includegraphics[width=\linewidth]{graficos/g6b_heatmap_sig.pdf}
  \caption{Mapa de calor do tempo de Sign (ML-DSA) e KeyGen
  (FN-DSA) por configuração de compilação. O padrão irregular
  do ML-DSA contrasta com o gradiente estruturado do FN-DSA.
  Valores normalizados por linha.}
  \label{fig:heatmap-sig}
\end{figure}
A Figura~\ref{fig:heatmap-sig} expõe padrões contrastantes
entre as duas famílias. Para o ML-DSA, o mapa de calor é
irregular: não há gradiente consistente entre configurações,
e células verdes e vermelhas alternam sem padrão direcional.
O ML-DSA-65, por exemplo, registra Sign de 88,5~$\mu$s com
\texttt{-O2\_native\_lto} (speedup de 1,09$\times$) e
104,9~$\mu$s com \texttt{-O1\_native\_lto} (0,92$\times$),
mas ambas as diferenças estão dentro da margem de dispersão
natural do mecanismo de \textit{Fiat-Shamir with Aborts},
cujo coeficiente de variação supera 50\%. A variabilidade
estocástica do algoritmo domina o efeito das flags, tornando
qualquer configuração equivalente na prática.
O FN-DSA exibe comportamento oposto: o padrão é estruturado
e consistente entre variantes. Para o KeyGen, a configuração
\texttt{-O2~-march=native~-flto} produz os menores tempos
em ambas as variantes, com 4.880,6~$\mu$s no FN-DSA-512
(speedup de 1,007$\times$) e 14.364,2~$\mu$s no FN-DSA-1024
(1,028$\times$). A configuração \texttt{-O2~-march=native}
sem LTO apresenta, de forma contraintuitiva, os piores
resultados: 5.459,7~$\mu$s no FN-DSA-512 (0,90$\times$),
sugerindo que \texttt{-march=native} isolado pode introduzir
pressão de registradores que prejudica a fase de amostragem
gaussiana, efeito mitigado pelo LTO ao otimizar os limites
entre módulos.
\subsubsection{SLH-DSA}
\begin{figure}[!htbp]
  \centering
  \includegraphics[width=\linewidth]{graficos/g6c_heatmap_slh.pdf}
  \caption{Mapa de calor do tempo de Sign do SLH-DSA-SHAKE
  por configuração de compilação. A configuração
  \texttt{-O2~-march=native~-flto} é consistentemente a mais
  rápida; \texttt{-O2~-flto} é consistentemente a mais lenta.
  Valores em ms, normalizados por linha.}
  \label{fig:heatmap-slh}
\end{figure}
O mapa de calor da Figura~\ref{fig:heatmap-slh} exibe o
padrão mais estruturado e interpretável dentre todas as
famílias avaliadas. Para as seis variantes do SLH-DSA, a
configuração \texttt{-O2~-march=native~-flto} concentra
consistentemente as células mais verdes, enquanto
\texttt{-O2~-flto} (sem \texttt{native}) concentra as
mais vermelhas. Essa oposição é visível em todas as linhas
do heatmap.
O SLH-DSA-SHAKE-128s registra Sign de 260.056~$\mu$s com
\texttt{-O2\_native\_lto} e 281.025~$\mu$s com
\texttt{-O2\_lto}, speedups de 1,001$\times$ e 0,927$\times$
respectivamente. O SLH-DSA-SHAKE-256s apresenta o maior
ganho absoluto: 405.773~$\mu$s com \texttt{-O2\_native\_lto}
frente a 440.057~$\mu$s com \texttt{-O2\_lto}, diferença
de 34,3~ms por operação de assinatura.
A explicação reside na estrutura computacional do algoritmo:
o SLH-DSA realiza centenas de avaliações de SHAKE-256 por
assinatura, operação que se beneficia diretamente das
extensões SHA e AVX2 habilitadas por \texttt{-march=native}.
O LTO isolado não acessa essas instruções por não ter
informação da arquitetura alvo, introduzindo apenas overhead
de linkedição sem ganho compensatório. Isso explica o fato
de \texttt{-O2\_lto} produzir tempos superiores ao próprio
\texttt{-O0} para o SLH-DSA.
Vale destacar ainda um padrão transversal: a configuração
\texttt{-O3} frequentemente regride em relação ao
\texttt{-O2} da mesma família de flags. O
SLH-DSA-SHAKE-128s, por exemplo, registra Sign de
260.550~$\mu$s com \texttt{-O3} e 260.056~$\mu$s com
\texttt{-O2\_native\_lto}. Esse comportamento é atribuído
ao aumento do tamanho do código gerado pelas heurísticas
agressivas do \texttt{-O3}: o desenrolamento excessivo de
loops e o inlining profundo produzem binários maiores que
exercem maior pressão sobre o cache de instruções (L1i),
cujo impacto é especialmente pronunciado em funções
altamente iterativas como a construção das árvores de Merkle.
Para o SLH-DSA, o ponto de equilíbrio entre otimização e
pressão de cache encontra-se em \texttt{-O2}, não em
\texttt{-O3}.
Os resultados desta seção evidenciam três regimes distintos
de sensibilidade às flags de compilação. O ML-KEM é
insensível porque seu núcleo crítico já está otimizado em
assembly dentro da liboqs. O ML-DSA é governado por
variabilidade estocástica, tornando qualquer conclusão sobre
flags estatisticamente frágil. O SLH-DSA e o FN-DSA
respondem de forma estruturada, com
\texttt{-O2~-march=native~-flto} como configuração
dominante, especialmente em operações que envolvem hashing
intensivo ou amostragem sobre redes de reticulados.
% ------------------------------------------------------------
% 5.4 — Escalabilidade com o Tamanho da Mensagem
% ------------------------------------------------------------
\subsection{Escalabilidade com o Tamanho da Mensagem}
Esta seção investiga como os tempos de assinatura e
verificação variam em função do tamanho da mensagem, de
1~B a 1~MiB em potências de dois (21 pontos), com a
configuração \texttt{-O0} como referência. A análise não
se aplica aos KEMs, cujas operações possuem entradas de
tamanho fixo. As Figuras~\ref{fig:escala-sign-lattice},
\ref{fig:escala-sign-slh-f} e \ref{fig:escala-sign-slh-s}
cobrem as curvas de assinatura; as
Figuras~\ref{fig:escala-verify-lattice} e
\ref{fig:escala-verify-slh} cobrem as de verificação.
\subsubsection{Assinatura: ML-DSA e FN-DSA}
\begin{figure}[!htbp]
  \centering
  \includegraphics[width=\linewidth]{graficos/g3a_escala_sign_lattice.pdf}
  \caption{Tempo de assinatura em função do tamanho da
  mensagem para ML-DSA e FN-DSA (\texttt{-O0}). Todas as
  curvas exibem crescimento linear com inclinação
  $\approx$2,1~$\mu$s/KiB a partir de 1~KiB, dominada
  pelo custo do SHAKE-256.}
  \label{fig:escala-sign-lattice}
\end{figure}
A Figura~\ref{fig:escala-sign-lattice} mostra que todas as
variantes do ML-DSA e do FN-DSA exibem o mesmo padrão:
tempo praticamente constante de 1~B até aproximadamente
1~KiB, seguido de crescimento linear com o tamanho da
mensagem. A inclinação dessa fase linear é praticamente
idêntica para todos os algoritmos, entre 2,08 e
2,16~$\mu$s/KiB, o que evidencia que o gargalo de
escalabilidade não é o núcleo criptográfico de cada
esquema, mas a função SHAKE-256 aplicada à mensagem antes
da assinatura, cuja taxa de processamento no AMD Ryzen 7
7435HS corresponde a aproximadamente 476~MB/s.
Embora a inclinação seja a mesma, os fatores de crescimento
de 1~B a 1~MiB diferem substancialmente: o ML-DSA-44, com
tempo base de 56,6~$\mu$s, cresce 38,7$\times$ até
2.189,5~$\mu$s; o FN-DSA-1024, com base de 339,2~$\mu$s,
cresce apenas 7,3$\times$ até 2.477,7~$\mu$s. Em valores
absolutos, todos os algoritmos convergem para a faixa de
2.189 a 2.478~$\mu$s no 1~MiB, confirmando que o hash
domina o tempo total para mensagens grandes.
Uma implicação prática direta desse resultado é que o custo
de assinar mensagens de diferentes tamanhos com ML-DSA ou
FN-DSA é fortemente dependente do volume: assinar um
documento de 1~MiB com ML-DSA-44 custa 38,7$\times$ mais
que assinar uma mensagem de 1~B. Para aplicações que
processam documentos grandes com frequência, como sistemas
de assinatura de contratos ou registros eletrônicos, esse
crescimento linear pode tornar-se um gargalo, e a análise
de desempenho deve ser conduzida com mensagens
representativas do caso de uso real, não com mensagens
mínimas.
\subsubsection{Assinatura: SLH-DSA}
\begin{figure}[!htbp]
  \centering
  \includegraphics[width=\linewidth]{graficos/g3b_escala_sign_slh_f.pdf}
  \caption{Tempo de assinatura em função do tamanho da
  mensagem para as variantes \textit{f} do SLH-DSA
  (\texttt{-O0}). As curvas são aproximadamente constantes;
  o crescimento de 1~B a 1~MiB não ultrapassa 33\%.}
  \label{fig:escala-sign-slh-f}
\end{figure}
\begin{figure}[!htbp]
  \centering
  \includegraphics[width=\linewidth]{graficos/g3c_escala_sign_slh_s.pdf}
  \caption{Tempo de assinatura em função do tamanho da
  mensagem para as variantes \textit{s} do SLH-DSA
  (\texttt{-O0}). O crescimento de 1~B a 1~MiB é inferior
  a 3,2\% em todas as variantes.}
  \label{fig:escala-sign-slh-s}
\end{figure}
As Figuras~\ref{fig:escala-sign-slh-f} e
\ref{fig:escala-sign-slh-s} apresentam o resultado mais
expressivo desta análise: o tempo de assinatura do SLH-DSA
é praticamente insensível ao tamanho da mensagem.
Para as variantes \textit{f}
(Figura~\ref{fig:escala-sign-slh-f}), as curvas permanecem
essencialmente planas ao longo de toda a faixa avaliada,
com crescimento restrito às mensagens maiores. O
SLH-DSA-SHAKE-128f varia de 13.401~$\mu$s (1~B) a
17.801~$\mu$s (1~MiB), fator de 1,33$\times$. O
SLH-DSA-SHAKE-256f registra variação de 43.380~$\mu$s a
47.720~$\mu$s, fator de apenas 1,10$\times$. As inclinações
das variantes \textit{f}, entre 4,20 e 4,26~$\mu$s/KiB,
são cerca do dobro das observadas para ML-DSA e FN-DSA,
indicando que o SLH-DSA aplica internamente um hashing mais
extenso da mensagem antes de iniciar a construção das
árvores de Merkle.
Para as variantes \textit{s}
(Figura~\ref{fig:escala-sign-slh-s}), o comportamento
aproxima-se ainda mais do constante. O SLH-DSA-SHAKE-128s
varia de 260.269~$\mu$s a 268.518~$\mu$s, acréscimo de
apenas 3,2\% sobre um custo base de 260~ms. O
SLH-DSA-SHAKE-192s cresce 1,2\%, e o SLH-DSA-SHAKE-256s,
1,3\%. As inclinações entre 4,23 e 6,30~$\mu$s/KiB
representam fração desprezível do custo total: para o
SLH-DSA-192s, o hash de 1~MiB adiciona cerca de 4,3~ms
sobre um custo base de 460~ms, contribuição inferior a 1\%.
Esse comportamento é consequência direta da estrutura do
algoritmo: o custo dominante é a construção das árvores
de Merkle da estrutura FORS e da hiper-árvore
XMS\textsuperscript{MT}, fixo e independente da mensagem.
A mensagem é hashada uma única vez antes de alimentar o
processo de assinatura, representando custo marginal no
cômputo total.
A insensibilidade ao tamanho da mensagem tem uma implicação
prática relevante: o SLH-DSA oferece latência de assinatura
previsível independentemente do volume de dados. Para cenários
de assinatura de documentos grandes, como PDFs ou imagens de
firmware, o SLH-DSA-SHAKE-128f registra tempos entre
13,4~ms e 17,8~ms de 1~B a 1~MiB, variação de apenas 4,4~ms.
O ML-DSA-44 é mais rápido para mensagens pequenas, mas o
SLH-DSA permanece mais lento em termos absolutos para qualquer
tamanho avaliado. Em contrapartida, a \textit{estabilidade}
do seu tempo de assinatura pode ser vantajosa em sistemas de
tempo real que demandam garantias de latência máxima.
\subsubsection{Verificação: ML-DSA e FN-DSA}
\begin{figure}[!htbp]
  \centering
  \includegraphics[width=\linewidth]{graficos/g5a_escala_verify_lattice.pdf}
  \caption{Tempo de verificação em função do tamanho da
  mensagem para ML-DSA e FN-DSA (\texttt{-O0}). A inclinação
  de $\approx$2,1~$\mu$s/KiB é idêntica à da assinatura,
  confirmando que o gargalo é o SHAKE-256.}
  \label{fig:escala-verify-lattice}
\end{figure}
A Figura~\ref{fig:escala-verify-lattice} mostra que a
verificação segue o mesmo padrão da assinatura: constante
até aproximadamente 1~KiB, linear a partir daí, com
inclinação de 2,08 a 2,16~$\mu$s/KiB. Os fatores de
crescimento de 1~B a 1~MiB são mais elevados que os da
assinatura porque os tempos base de verificação são
menores: o ML-DSA-44 cresce 100,5$\times$ (de 21,4~$\mu$s
a 2.152~$\mu$s) e o FN-DSA-512 cresce 65,1$\times$ (de
34,0~$\mu$s a 2.211~$\mu$s). As cinco curvas convergem
para valores próximos em 1~MiB (entre 2.152 e 2.250~$\mu$s),
reforçando que o SHAKE-256 é o gargalo universal da fase
de verificação para mensagens grandes.
A uniformidade da inclinação de $\approx$2,1~$\mu$s/KiB
entre todos os algoritmos na verificação é um resultado que
transcende a comparação de desempenho: indica que, para
mensagens grandes, o custo de verificação de uma assinatura
digital baseada em SHAKE-256 é determinado essencialmente
pelo hardware e pelo tamanho da mensagem, sendo
independente do esquema criptográfico adotado. Para o
dimensionamento de sistemas de verificação em lote, isso
significa que verificar 1.000 assinaturas sobre documentos
de 1~MiB custará aproximadamente o mesmo tempo de hash para
ML-DSA-44 ou FN-DSA-512.
\subsubsection{Verificação: SLH-DSA}
\begin{figure}[!htbp]
  \centering
  \includegraphics[width=\linewidth]{graficos/g5b_escala_verify_slh.pdf}
  \caption{Tempo de verificação em função do tamanho da
  mensagem para todas as variantes do SLH-DSA (\texttt{-O0}).
  Ao contrário da assinatura, a verificação cresce linearmente
  em todas as variantes com inclinação $\approx$2,1~$\mu$s/KiB.
  Linhas sólidas representam a variante \textit{s}; tracejadas,
  a variante \textit{f}.}
  \label{fig:escala-verify-slh}
\end{figure}
A Figura~\ref{fig:escala-verify-slh} reúne as curvas de
verificação das seis variantes do SLH-DSA. Ao contrário
da assinatura, a verificação cresce linearmente com a
mensagem em todas as variantes, com inclinação entre 2,12
e 2,14~$\mu$s/KiB, valor que coincide com o observado para
ML-DSA e FN-DSA. Isso confirma que, na fase de verificação,
o SLH-DSA percorre apenas um caminho da árvore de Merkle
(custo fixo e pequeno) e hasha a mensagem completa, operação
que passa a dominar o tempo para mensagens maiores.
As variantes \textit{f} (tracejadas) partem de tempos base
superiores aos das \textit{s} (sólidas) do mesmo nível:
o SLH-128f inicia em 898~$\mu$s, enquanto o SLH-128s inicia
em 331~$\mu$s. Ambas convergem para valores próximos em
1~MiB (3.077 e 2.517~$\mu$s, respectivamente). A redução
gradual desse intervalo ocorre porque a fração de tempo
dedicada ao hash da mensagem torna-se dominante em ambas
as variantes, diluindo a diferença nos tempos base.
Do ponto de vista prático, a verificação do SLH-DSA
apresenta um perfil favorável para aplicações de alta
frequência com mensagens grandes: o tempo base de
verificação para mensagens pequenas varia de 331~$\mu$s
a 1.309~$\mu$s conforme a variante, mas a inclinação de
2,1~$\mu$s/KiB é a mesma dos demais esquemas. Para mensagens
de 1~MiB, todas as variantes convergem para a faixa de
2,5 a 4,1~ms, apenas marginalmente superior aos 2,2~ms do
ML-DSA e FN-DSA. Em sistemas onde a verificação de
documentos grandes é a operação dominante, a penalidade de
verificação do SLH-DSA é pequena em termos absolutos, mesmo
que o custo de assinatura seja ordens de grandeza mais
elevado.
% ------------------------------------------------------------
% 5.5 — Consumo de Memória
% ------------------------------------------------------------
\subsection{Consumo de Memória}
\label{sec:res-memoria}


O consumo de memória foi avaliado por duas técnicas
complementares: \textit{stack painting} para medir o uso
de pilha e \texttt{valgrind --tool=massif} para medir o
pico de \textit{heap}. Ambas as medições foram realizadas
com compilação \texttt{-O0}, sem otimizações, para refletir
o comportamento real de alocação de cada algoritmo.
É importante notar que o pico de \textit{heap} capturado
pelo Valgrind mede o máximo de memória alocada \emph{durante}
a execução, incluindo alocações transientes que são liberadas
ao término da operação; o delta líquido de \textit{heap}
entre o início e o fim de cada operação é nulo, pois o
\textit{liboqs} não deixa alocações persistentes após o
retorno das funções.


\subsubsection{Uso de Pilha (\textit{Stack})}
\label{sec:res-stack}


As Tabelas~\ref{tab:stack_kem} e~\ref{tab:stack_sig}
apresentam o uso de \textit{stack} para os KEMs e para
os esquemas de assinatura, respectivamente.


\begin{table}[ht]
\centering
\caption{Uso de memória \textit{stack} por KEM e operação
(bytes). Medido via \textit{stack painting} com
\texttt{pthread} de 4\,MB e compilação \texttt{-O0}.}
\label{tab:stack_kem}
\resizebox{\columnwidth}{!}{%
\begin{tabular}{lrrr}
\toprule
\textbf{Algoritmo} & \textbf{Keygen} & \textbf{Encap} & \textbf{Decap} \\
\midrule
\multicolumn{4}{l}{\textit{ML-KEM (FIPS\,203)}} \\[2pt]
\quad ML-KEM-512  &  13.520 &  16.160 &  16.960 \\
\quad ML-KEM-768  &  17.648 &  20.784 &  21.904 \\
\quad ML-KEM-1024 &  22.768 &  26.448 &  28.048 \\
\midrule
\multicolumn{4}{l}{\textit{HQC}} \\[2pt]
\quad HQC-128 &  38.576 &  56.448 &  60.896 \\
\quad HQC-192 &  70.688 & 106.752 & 115.760 \\
\quad HQC-256 & 108.768 & 166.592 & 181.040 \\
\bottomrule
\end{tabular}}
\end{table}


\begin{table}[ht]
\centering
\caption{Uso de memória \textit{stack} por esquema de
assinatura e operação (bytes). Medido via \textit{stack
painting} com \texttt{pthread} de 4\,MB e compilação
\texttt{-O0}.}
\label{tab:stack_sig}
\resizebox{\columnwidth}{!}{%
\begin{tabular}{lrrr}
\toprule
\textbf{Algoritmo} & \textbf{Keygen} & \textbf{Sign} & \textbf{Verify} \\
\midrule
\multicolumn{4}{l}{\textit{ML-DSA (FIPS\,204)}} \\[2pt]
\quad ML-DSA-44 &  28.304 &  53.456 &  25.968 \\
\quad ML-DSA-65 &  33.456 &  83.312 &  29.072 \\
\quad ML-DSA-87 &  41.648 & 117.072 &  35.472 \\
\midrule
\multicolumn{4}{l}{\textit{FN-DSA (FIPS\,206)}} \\[2pt]
\quad FN-DSA-512  &  23.608 &  49.584 &   9.064 \\
\quad FN-DSA-1024 &  40.504 &  90.736 &  13.160 \\
\midrule
\multicolumn{4}{l}{\textit{SLH-DSA-SHAKE (FIPS\,205)}} \\[2pt]
\quad SLH-DSA-SHAKE-128s &  14.480 &  13.744 &   8.464 \\
\quad SLH-DSA-SHAKE-128f &  14.000 &  13.360 &   8.464 \\
\quad SLH-DSA-SHAKE-192s &  20.976 &  19.504 &   9.328 \\
\quad SLH-DSA-SHAKE-192f &  20.240 &  18.928 &   9.360 \\
\quad SLH-DSA-SHAKE-256s &  29.744 &  27.312 &  10.480 \\
\quad SLH-DSA-SHAKE-256f &  29.104 &  26.800 &  10.512 \\
\bottomrule
\end{tabular}}
\end{table}


O uso de \textit{stack} é a métrica que melhor diferencia
os algoritmos sob análise.
Entre os KEMs, o ML-KEM apresenta consumo reduzido e
previsível: a operação mais custosa, \textit{decap} do
ML-KEM-1024, consome apenas 28.048 bytes. O HQC, por outro
lado, exige consideravelmente mais pilha: o \textit{decap}
do HQC-256 chega a 181.040 bytes, cerca de 6,5 vezes o
valor equivalente do ML-KEM-1024. Esse comportamento é
esperado dado o tamanho dos vetores polinomiais e os buffers
intermediários que o HQC emprega internamente.


Entre os esquemas de assinatura, o ML-DSA é o mais exigente:
a operação \textit{sign} do ML-DSA-87 consome 117.072 bytes,
reflexo dos vetores de comprometimento e das matrizes de
desafio manipulados durante a geração da assinatura. O
FN-DSA destaca-se pelo baixo consumo de \textit{stack} na
verificação: 9.064 bytes para FN-DSA-512 e 13.160 bytes
para FN-DSA-1024.


O resultado mais surpreendente desta análise pertence ao
SLH-DSA-SHAKE: apesar de registrar o maior custo
computacional de assinatura de todo o conjunto avaliado
(até 460\,ms), seu uso de \textit{stack} é o menor entre
todos os esquemas de assinatura. A verificação das variantes
SLH-DSA-SHAKE-128s e -128f consome apenas 8.464 bytes,
valor inferior ao do próprio FN-DSA-512 (9.064 bytes) na
mesma operação. As operações de assinatura ficam entre
13.360 bytes (SLH-DSA-SHAKE-128f) e 27.312 bytes
(SLH-DSA-SHAKE-256s), intervalo comparável ao do ML-KEM-512.
Essa propriedade é consequência direta da arquitetura
baseada em funções hash: o SLH-DSA não emprega matrizes
nem vetores polinomiais de grande porte alocados em pilha;
toda a computação se realiza por chamadas iterativas de
SHAKE-256 com estado interno reduzido. Assim, o SLH-DSA
combina alto custo de CPU com uso mínimo de \textit{stack},
tornando-o uma opção viável em ambientes com restrições
severas de pilha que toleram latência elevada de assinatura.


\subsubsection{Pico de Heap (Valgrind \texttt{massif})}
\label{sec:res-heap}


As Tabelas~\ref{tab:heap_kem} e~\ref{tab:heap_sig}
apresentam o pico de uso de \textit{heap} medido pelo
\texttt{valgrind --tool=massif}. Para as variantes
SLH-DSA-SHAKE, os dados de \textit{heap} estão disponíveis,
pois foram coletados com o programa \texttt{massif\_single.c}
de forma independente do benchmark de \textit{stack}.


\begin{table}[ht]
\centering
\caption{Pico de \textit{heap} por KEM e operação (KB).
Medido com \texttt{valgrind --tool=massif}; valor de pico
de \texttt{mem\_heap\_B} para cada operação isolada.}
\label{tab:heap_kem}
\resizebox{\columnwidth}{!}{%
\begin{tabular}{lrrr}
\toprule
\textbf{Algoritmo} & \textbf{Keygen (KB)} & \textbf{Encap (KB)} & \textbf{Decap (KB)} \\
\midrule
\multicolumn{4}{l}{\textit{ML-KEM (FIPS\,203)}} \\[2pt]
\quad ML-KEM-512  & 204,97 & 204,97 & 204,97 \\
\quad ML-KEM-768  & 207,17 & 207,17 & 207,17 \\
\quad ML-KEM-1024 & 208,00 & 208,17 & 208,17 \\
\midrule
\multicolumn{4}{l}{\textit{HQC}} \\[2pt]
\quad HQC-128 & 211,07 & 211,07 & 211,07 \\
\quad HQC-192 & 219,74 & 219,74 & 219,74 \\
\quad HQC-256 & 230,60 & 230,60 & 230,60 \\
\bottomrule
\end{tabular}}
\end{table}


\begin{table}[ht]
\centering
\caption{Pico de \textit{heap} por esquema de assinatura
digital e operação (KB). Medido com
\texttt{valgrind --tool=massif}; valor de pico de
\texttt{mem\_heap\_B} para cada operação isolada.}
\label{tab:heap_sig}
\resizebox{\columnwidth}{!}{%
\begin{tabular}{lrrr}
\toprule
\textbf{Algoritmo} & \textbf{Keygen (KB)} & \textbf{Sign (KB)} & \textbf{Verify (KB)} \\
\midrule
\multicolumn{4}{l}{\textit{ML-DSA (FIPS\,204)}} \\[2pt]
\quad ML-DSA-44 & 207,92 & 208,90 & 208,90 \\
\quad ML-DSA-65 & 211,02 & 211,02 & 211,83 \\
\quad ML-DSA-87 & 214,37 & 214,59 & 214,59 \\
\midrule
\multicolumn{4}{l}{\textit{FN-DSA (FIPS\,206)}} \\[2pt]
\quad FN-DSA-512  & 204,64 & 204,81 & 204,81 \\
\quad FN-DSA-1024 & 207,16 & 207,38 & 207,38 \\
\midrule
\multicolumn{4}{l}{\textit{SLH-DSA-SHAKE (FIPS\,205)}} \\[2pt]
\quad SLH-DSA-SHAKE-128s & 210,31 & 210,31 & 210,31 \\
\quad SLH-DSA-SHAKE-128f & 219,32 & 219,32 & 219,32 \\
\quad SLH-DSA-SHAKE-192s & 218,52 & 218,52 & 218,52 \\
\quad SLH-DSA-SHAKE-192f & 237,51 & 237,51 & 237,51 \\
\quad SLH-DSA-SHAKE-256s & 231,82 & 231,82 & 231,82 \\
\quad SLH-DSA-SHAKE-256f & 251,42 & 251,42 & 251,42 \\
\bottomrule
\end{tabular}}
\end{table}


Ao contrário do \textit{stack}, o pico de \textit{heap}
não constitui um fator diferenciador expressivo entre os
algoritmos. Os valores observados situam-se em uma faixa
estreita, de aproximadamente 204,64\,KB (FN-DSA-512) a
251,42\,KB (SLH-DSA-SHAKE-256f), uma variação total
inferior a 47\,KB em todo o conjunto avaliado.


Essa uniformidade é explicada pela arquitetura interna do
\texttt{liboqs}: a biblioteca realiza uma alocação de
contexto durante a inicialização do objeto de algoritmo
(\texttt{OQS\_KEM\_new} / \texttt{OQS\_SIG\_new}), que
representa o grosso do \textit{heap} medido pelo Valgrind,
independentemente do algoritmo selecionado. As alocações
transitórias da operação em si são adicionadas durante a
execução e liberadas ao término, resultando em variações
pequenas entre operações do mesmo algoritmo.


O crescimento do pico entre variantes de maior nível de
segurança reflete principalmente o tamanho dos buffers
internos: o SLH-DSA-SHAKE-256f, com parâmetros mais
conservadores e assinaturas maiores, apresenta o maior
pico (251,42\,KB), enquanto o FN-DSA-512 registra o menor
(204,64\,KB). Para efeitos de seleção de algoritmo em
sistemas com restrição de memória dinâmica, o \textit{heap}
não é o fator limitante: todos os algoritmos operam dentro
de uma faixa compatível com sistemas de propósito geral,
e o \textit{stack} é o critério relevante para ambientes
com restrições severas de memória.
% ============================================================
% SEÇÃO 6 — Conclusão
% ============================================================
\section{Conclusão}
Este trabalho apresentou uma análise experimental de
desempenho dos algoritmos finalistas do processo de
padronização do NIST para criptografia pós-quântica,
avaliando ML-KEM, HQC, ML-DSA, FN-DSA e SLH-DSA sob
cinco dimensões: tamanho dos artefatos criptográficos,
desempenho temporal das operações fundamentais, sensibilidade
às configurações de compilação, escalabilidade com o
tamanho da mensagem e consumo de memória. Os experimentos
foram conduzidos com a biblioteca liboqs sobre hardware
x86-64 com suporte a instruções AVX2, SHA e FMA3, cobrindo
16 configurações de compilação e todas as variantes
paramétricas de cada algoritmo nos níveis de segurança
NIST 1, 3 e 5.
Os resultados evidenciam que nenhum algoritmo domina todas
as dimensões simultaneamente, e que a escolha adequada
depende do equilíbrio entre as propriedades relevantes para
cada aplicação. No grupo dos KEMs, o ML-KEM supera o HQC
em desempenho por um fator de 200$\times$ a 1.503$\times$
dependendo da operação e do nível de segurança, beneficiado
por implementações com instruções SIMD e pela eficiência da
Transformada Número-Teórica sobre módulos de reticulados.
O HQC, apesar do custo computacional elevado, oferece
diversificação da base matemática ao fundamentar-se em
teoria de códigos corretores de erros, propriedade relevante
em cenários que demandam resistência a múltiplas classes de
ataques.
No grupo das assinaturas digitais, o ML-DSA apresenta o
melhor equilíbrio geral: as três operações fundamentais
operam na mesma ordem de grandeza, e o algoritmo é
competitivo em todos os níveis de segurança avaliados. Seu
principal ponto de atenção é o alto coeficiente de variação
do Sign (51\% a 62\%), decorrente do mecanismo de
\textit{Fiat-Shamir with Aborts}, que torna a média
estatisticamente menos representativa que a mediana para
dimensionamento de sistemas. O FN-DSA ocupa uma posição
distinta: o custo de geração de chave é 28$\times$ superior
ao de assinatura no nível 1, tornando-o desfavorável em
cenários de chave efêmera, mas competitivo quando o par de
chaves é reutilizado ao longo de múltiplas operações. O
SLH-DSA concentra o maior custo computacional de assinatura
(até 460~ms na variante \textit{s}, nível 3), compensado
por duas propriedades estruturais: comportamento
determinístico e independência quase total do tempo de
assinatura em relação ao tamanho da mensagem, com crescimento
inferior a 3,2\% de 1~B a 1~MiB.


O consumo de memória revelou que o \textit{stack} é a
dimensão mais discriminante entre os algoritmos nessa
métrica. O HQC exige até 6,5 vezes mais pilha que o ML-KEM
no nível de segurança equivalente, reflexo de seu design
baseado em vetores densos. O resultado mais inesperado
desta dimensão pertence ao SLH-DSA-SHAKE: apesar de
registrar o maior custo computacional de assinatura, seu
uso de \textit{stack} é o mais baixo de todo o conjunto —
a verificação do SLH-DSA-SHAKE-128s consome apenas
8.464 bytes, superando inclusive o FN-DSA-512 (9.064 bytes)
nessa métrica. Essa propriedade decorre de sua arquitetura
iterativa baseada em funções hash, sem alocação de matrizes
polinomiais em pilha. O pico de \textit{heap}, por sua vez,
é dominado pelo \textit{overhead} de inicialização do
\texttt{liboqs} e varia menos de 47\,KB entre todos os
algoritmos avaliados, não constituindo critério
diferenciador para seleção em sistemas de propósito geral.


A análise das configurações de compilação revelou três
regimes distintos de sensibilidade. O ML-KEM é insensível
às flags do código cliente, pois seu núcleo criptográfico
já está otimizado em assembly dentro da liboqs, e as
variações observadas são estatisticamente indistinguíveis
do ruído de medição. O ML-DSA é governado pela
variabilidade estocástica do algoritmo, que domina qualquer
efeito de compilação. O FN-DSA e o SLH-DSA respondem de
forma estruturada, com a combinação
\texttt{-O2~-march=native~-flto} produzindo consistentemente
os melhores tempos ao habilitar instruções SHA e AVX2 da
arquitetura. Um achado transversal é que o nível \texttt{-O3}
frequentemente regride em relação ao \texttt{-O2} por
introduzir pressão excessiva sobre o cache de instruções,
sugerindo que o ponto de equilíbrio entre otimização e
localidade de cache encontra-se em \texttt{-O2} para
algoritmos com funções altamente iterativas.
A análise de escalabilidade identificou um gargalo universal:
a função SHAKE-256 aplicada à mensagem impõe uma inclinação
de $\approx$2,1~$\mu$s/KiB em todas as operações de
assinatura e verificação, para todos os algoritmos avaliados.
Esse resultado indica que, para mensagens grandes, o custo
de assinatura e verificação é determinado primariamente pelo
hardware e pelo tamanho da mensagem, e não pelo esquema
criptográfico adotado. A exceção é o Sign do SLH-DSA, cujo
custo base das árvores de Merkle domina completamente o hash
da mensagem, tornando-o praticamente constante ao longo de
toda a faixa avaliada.
Dois achados merecem destaque como contribuições específicas
deste trabalho. O primeiro é a quantificação do viés
introduzido pela amostragem por rejeição no ML-DSA: a média
do Sign supera a mediana em 27,7\% no ML-DSA-44, diferença
que decresce para 18,1\% no ML-DSA-87 à medida que os
parâmetros do nível de segurança mais alto reduzem a taxa
de rejeição do protocolo. Essa assimetria implica que
sistemas dimensionados com base na média aritmética
subestimarão o tempo real de execução em até 28\% para
um conjunto típico de requisições, e que a mediana é a
estatística adequada para planejamento de capacidade em
aplicações baseadas em ML-DSA. O segundo achado é a
inclinação universal de $\approx$2,1~$\mu$s/KiB imposta
pelo SHAKE-256 em todas as operações de assinatura e
verificação de todos os algoritmos avaliados, indicando
que o gargalo de escalabilidade para mensagens grandes
é determinado pelo hardware, não pelo esquema criptográfico.
Nenhum dos trabalhos correlatos revisados reporta qualquer
um desses dois resultados.
O trabalho apresenta duas limitações principais.
Os experimentos foram conduzidos em uma única plataforma
de hardware sob WSL, e os resultados podem diferir em
ambientes de execução distintos, como servidores bare metal,
processadores ARM ou plataformas embarcadas com suporte
limitado a extensões SIMD.
Por fim, a análise das configurações de compilação avalia
exclusivamente o desempenho temporal, sem verificar a
preservação das propriedades de segurança de cada algoritmo
sob diferentes níveis de otimização. Implementações
criptográficas seguras frequentemente dependem de execução
em tempo constante (\textit{constant-time}) para resistir
a ataques de canal lateral por temporização. Otimizações
agressivas do compilador, como vetorização automática e
reorganização de instruções presentes nos níveis
\texttt{-O2} e \texttt{-O3}, podem alterar os padrões de
execução de forma a comprometer essas garantias. Assim,
os ganhos de desempenho identificados na Seção~5.3 não
devem ser interpretados como recomendação de adoção direta
dessas configurações sem validação independente das
propriedades de segurança, especialmente em contextos onde
a resistência a ataques de canal lateral é requisito.
Os resultados obtidos respondem à questão central proposta
na introdução: as diferenças de desempenho entre os
algoritmos finalistas do NIST PQC são substanciais em todas
as dimensões avaliadas, e nenhum esquema domina os demais
em todas as métricas simultaneamente. Essa diversidade de
perfis não é uma limitação do processo de padronização;
pelo contrário, reflete a intenção deliberada do NIST de
oferecer alternativas adequadas a contextos distintos,
desde aplicações de baixa latência até cenários que exigem
assinaturas compactas, comportamento determinístico ou
diversificação da base matemática. A transição da
criptografia clássica para o paradigma pós-quântico
demandará escolhas informadas por dados experimentais,
e este trabalho contribui com uma base comparativa
reprodutível para subsidiar essas decisões. Os resultados
foram exportados em formato CSV padronizado, com metodologia
documentada e verificação de corretude por iteração,
permitindo reprodução independente e extensão para novos
algoritmos, plataformas ou métricas em estudos futuros.
\subsection{Trabalhos Futuros}
Os resultados e limitações identificados neste trabalho
apontam para diversas direções de pesquisa que podem
aprofundar ou complementar a análise realizada.
\textbf{Novos algoritmos do processo de padronização do NIST.}
O processo de padronização conduzido pelo NIST ainda não
está encerrado: rodadas adicionais avaliaram candidatos
suplementares, e novos esquemas podem ser incorporados ao
conjunto de algoritmos padronizados. Uma extensão natural
deste trabalho é a inclusão desses candidatos na mesma
infraestrutura de benchmark, permitindo comparação direta
com os algoritmos aqui avaliados sob as mesmas condições
experimentais.
\textbf{Validação das propriedades de segurança sob diferentes
flags de compilação.}
Este trabalho avaliou exclusivamente o impacto das flags
sobre o desempenho temporal, sem verificar se as
propriedades de segurança, em especial a execução em tempo
constante (\textit{constant-time}), são preservadas em cada
configuração. Trabalhos futuros podem empregar ferramentas
de análise estática e dinâmica de temporização, como
\cite{Reparaz:2017} e
\cite{Almeida:2016}, para identificar
configurações de compilação que introduzem dependência de
dados no tempo de execução, comprometendo a resistência a
ataques de canal lateral.
\textbf{Expansão das configurações de compilação.}
As 16 configurações avaliadas combinam quatro níveis de
otimização com \texttt{-march=native} e \texttt{-flto}.
Estudos futuros podem ampliar esse espaço com
\textit{Profile-Guided Optimization} (PGO), que utiliza
dados de execução real para guiar as decisões do compilador,
com a \textit{polyhedral model optimization}
(\texttt{-floop-nest-optimize}), e com comparações entre
compiladores distintos, como GCC e Clang/LLVM, cujas
heurísticas de vetorização diferem significativamente para
código com padrões de acesso irregular como NTT e hashing
iterativo.
\textbf{Avaliação em plataformas embarcadas e IoT.}
Os experimentos foram realizados sobre um processador
x86-64 com suporte completo a instruções AVX2 e SHA. O
desempenho relativo entre os algoritmos pode diferir
substancialmente em arquiteturas ARM Cortex-M ou RISC-V
sem extensões SIMD, cenário relevante para dispositivos
IoT e sistemas embarcados que constituem um dos principais
casos de uso da criptografia pós-quântica. Estudos em
plataformas como Raspberry Pi, STM32 e ESP32 permitiriam
avaliar a viabilidade prática dos algoritmos em ambientes
com restrições severas de processamento e energia.
\textbf{Integração com protocolos de comunicação reais.}
Este trabalho avaliou os algoritmos em isolamento. Uma
extensão relevante é a medição do impacto sobre protocolos
reais, como TLS~1.3 e SSH, onde o custo de cada operação
criptográfica é distribuído ao longo do handshake e de
sessões de longa duração. Estudos como o de
\cite{Sosnowski:2023} demonstraram que o desempenho de
PQC em contexto de protocolo pode diferir das medições
isoladas, e a análise combinada de KEMs e esquemas de
assinatura no mesmo handshake permitiria avaliar o custo
total da migração pós-quântica.
\textbf{Benchmarks com modo híbrido clássico e pós-quântico.}
A estratégia de migração recomendada pelo NIST e por
organismos como o IETF prevê um período de uso híbrido, em
que algoritmos clássicos (como X25519 e ECDSA) operam em
paralelo com os seus substitutos pós-quânticos. Avaliar o
overhead combinado de pares híbridos (por exemplo,
X25519~+~ML-KEM e Ed25519~+~ML-DSA) é uma extensão
direta deste trabalho com impacto prático imediato para
engenheiros e arquitetos de sistemas de segurança em
processo de migração.
\bibliographystyle{apalike-sol}
\bibliography{refs}
\end{document}
